% Estilo exposición Beamer — Grupo 9.2 (16:9, español, Helvetica)
\usepackage[utf8]{inputenc}
\usepackage[T1]{fontenc}
\usepackage[spanish,es-tabla]{babel}
\usepackage{graphicx}
\usepackage{adjustbox}
\usepackage{array}
\usepackage{colortbl}
\usepackage{booktabs}
\usepackage{ragged2e}
\usepackage{tabularx}
\usepackage{tcolorbox}
\tcbuselibrary{skins}
\usepackage{tikz}
\usepackage{xurl}

\usetheme{Madrid}
\usecolortheme{whale}
\usefonttheme{professionalfonts}
\usepackage[scaled=0.92]{helvet}
\renewcommand{\familydefault}{\sfdefault}

\definecolor{uninavy}{RGB}{27,54,93}
\definecolor{uniaccent}{RGB}{232,163,23}
\definecolor{unimuted}{RGB}{99,110,114}
\definecolor{unidanger}{RGB}{192,57,43}
\definecolor{tablehead}{RGB}{27,54,93}
\definecolor{tablealt}{RGB}{238,242,247}
\definecolor{tablerule}{RGB}{206,214,224}
\definecolor{tableframe}{RGB}{27,54,93}

\setbeamercolor{structure}{fg=uninavy}
\setbeamercolor{frametitle}{bg=uninavy,fg=white}
\setbeamercolor{title}{fg=white,bg=uninavy}
\setbeamercolor{block title}{bg=uninavy,fg=white}
\setbeamercolor{block body}{bg=uninavy!5,fg=black}
\setbeamercolor{block title alerted}{bg=unidanger,fg=white}
\setbeamercolor{block body alerted}{bg=unidanger!8,fg=black}
\setbeamercolor{block title example}{bg=uniaccent,fg=uninavy}
\setbeamercolor{block body example}{bg=uniaccent!14,fg=black}
\setbeamercolor{footline}{fg=unimuted,bg=white}

\setbeamertemplate{navigation symbols}{}
\setbeamertemplate{blocks}[rounded][shadow=false]
\setbeamersize{text margin left=11mm, text margin right=11mm}

\setbeamertemplate{footline}{%
  \leavevmode%
  \hbox{%
    \begin{beamercolorbox}[wd=.72\paperwidth,ht=2.5ex,dp=1.2ex,leftskip=1em]{footline}%
      \grupo{} · Plan de Tesis · UNI · Jul 2026%
    \end{beamercolorbox}%
    \begin{beamercolorbox}[wd=.28\paperwidth,ht=2.5ex,dp=1.2ex,rightskip=1em]{footline}%
      \insertframenumber{} / \inserttotalframenumber%
    \end{beamercolorbox}%
  }%
}

\setbeamertemplate{itemize item}{\color{uniaccent}\small$\blacktriangleright$}
\setbeamertemplate{itemize subitem}{\color{uninavy!70}\scriptsize$\circ$}
\setbeamerfont{normal text}{size=\large}
\setbeamerfont{item}{size=\large}
\setbeamerfont{frametitle}{size=\large,series=\bfseries}
\setbeamerfont{block title}{size=\normalsize,series=\bfseries}
\setbeamerfont{block body}{size=\large}

\graphicspath{%
  {ppt-images/}%
  {../../figuras/cap01/}%
  {../../figuras/cap02/}%
  {../../figuras/cap03/}%
  {../../figuras/modelo-alineacion/}%
}

\newcommand{\figplan}[2][]{%
  \vspace{-0.15em}%
  \centering
  \adjustbox{max width=0.94\linewidth, max height=0.68\textheight, center}{%
    \includegraphics[#1]{#2}%
  }%
}

\newcommand{\figplanlarge}[2][]{%
  \vspace{-0.35em}%
  \centering
  \adjustbox{max width=0.98\linewidth, max height=0.80\textheight, center}{%
    \includegraphics[#1]{#2}%
  }%
}

% Figura ontología (slide 4 · no oral) — máximo espacio vertical
\newcommand{\figplanontology}[2][]{%
  \vspace{-0.45em}%
  \centering
  \adjustbox{max width=\linewidth, max height=0.83\textheight, center}{%
    \includegraphics[#1]{#2}%
  }%
}

\newcommand{\figplancompact}[2][]{%
  \vspace{-0.15em}%
  \centering
  \adjustbox{max width=0.92\linewidth, max height=0.50\textheight, center}{%
    \includegraphics[#1]{#2}%
  }%
}

\newcommand{\figppttwo}[2]{%
  \vspace{-0.2em}%
  \centering
  \adjustbox{max width=0.96\linewidth, max height=0.34\textheight, center}{%
    \includegraphics{#1}%
  }\\[0.12em]
  \adjustbox{max width=0.96\linewidth, max height=0.34\textheight, center}{%
    \includegraphics{#2}%
  }%
}

\newcommand{\figcaption}[1]{%
  \vspace{0.2em}%
  \parbox{0.94\linewidth}{\centering\large\bfseries\textcolor{uninavy}{#1}}%
}

% --- Tablas ---
\newcolumntype{L}[1]{>{\RaggedRight\arraybackslash}p{#1}}
\newcolumntype{Y}{>{\RaggedRight\arraybackslash}X}
\newcommand{\tabhead}[1]{\color{white}\textbf{#1}}
\newcommand{\tabphase}[1]{\textcolor{uninavy}{\textbf{#1}}}

\newcommand{\plantablesetup}{%
  \setlength{\tabcolsep}{5.5pt}%
  \renewcommand{\arraystretch}{1.14}%
  \arrayrulecolor{tablerule}%
}

% Marco ancho completo · esquinas redondeadas · franja superior dorada
\newcommand{\beamertable}[1]{%
  \par\vspace{0.1em}%
  \begin{tcolorbox}[
    enhanced,
    arc=4pt,
    boxrule=0.55pt,
    colframe=tableframe,
    colback=white,
    width=\linewidth,
    boxsep=0pt,
    left=7pt, right=7pt, top=5pt, bottom=5pt,
    borderline north={1.05pt}{0pt}{uniaccent},
  ]
  {\fontsize{13.5}{16.2}\selectfont\plantablesetup #1}
  \end{tcolorbox}
  \par\vspace{0.2em}%
}

\newcommand{\beamertablecompact}[1]{%
  \par\vspace{0.03em}%
  \begin{tcolorbox}[
    enhanced,
    arc=3pt,
    boxrule=0.55pt,
    colframe=tableframe,
    colback=white,
    width=\linewidth,
    boxsep=0pt,
    left=4pt, right=4pt, top=3pt, bottom=3pt,
    borderline north={1.05pt}{0pt}{uniaccent},
  ]
  {%
    \fontsize{7.75}{9.4}\selectfont
    \plantablesetup
    \setlength{\tabcolsep}{3pt}%
    \renewcommand{\arraystretch}{1.0}%
    #1%
  }
  \end{tcolorbox}
  \par\vspace{0.1em}%
}

\newcommand{\tablesectiontitle}[1]{%
  \par\vspace{0.22em}%
  \noindent
  {\color{uniaccent}\rule{0.045\linewidth}{2pt}\hspace{0.35em}%
   \small\bfseries\textcolor{uninavy}{#1}}%
  \par\vspace{0.08em}%
}

% Bloque recordatorio (acento dorado)
\newenvironment{planreminder}[1]{%
  \begin{exampleblock}{#1}%
}{%
  \end{exampleblock}%
}

% Lista en panel suave (sin título)
\newenvironment{planpanel}{%
  \begin{beamercolorbox}[wd=0.98\linewidth,sep=0.55em,rounded=true,shadow=false]{block body}%
}{%
  \end{beamercolorbox}%
}

\newcommand{\keyterm}[1]{\textbf{\textcolor{uninavy}{#1}}}

% Punto rojo esquina superior derecha — slide no oral
\newcommand{\marcadorNoExponer}{%
  \begin{tikzpicture}[remember picture,overlay]
    \fill[unidanger] ([xshift=-0.65cm,yshift=-0.55cm]current page.north east) circle (6pt);
  \end{tikzpicture}%
}
